\scanpage{12}
\mspage{19}
% One continuous cancellation, made with a frame and diagonal strokes, reaches
% from (4.16) to the close of the paragraph before “Alors l'image”.
\[
\deleted{\handwritten{H^1(\unsure{\underline h|h|\underline k|k|y},\mathrm{Pic}_{X/h})=0}\quad
\deleted{\mathrm{Br}(X_{\handwritten{\unsure{y}}})=0,\quad H^{\handwritten3}(y,\underline G_m)=0,}\quad
\handwritten{\mathrm{Ker}(H^3(y,G_m)\longrightarrow H^3(X_y,G_m))}}
\tag*{\deleted{(4.1\deleted{\unsure{4}}\handwritten{6})}}
\]
\deleted{(Pour le voir, utiliser \deleted{\unsure{un|l'} homomorphisme de}\inserted{\handwritten{l'homomorphisme}} suite\handwritten{s} exacte\handwritten{s} du type (4.5) relatives à $\bar f:\bar X\handwritten{\longrightarrow}\bar Y$ \deleted{\unsure{et à $f_y:X_y\longrightarrow y$, qui donne,}}\inserted{\handwritten{et à $f_y:X_y\longrightarrow y$, qui s'écrit,}} \handwritten{compte tenu de (3.4) b) \unsure{:|;} et de (\phantom{000}) :}}
\diagram{\begin{tikzcd}[ampersand replacement=\&,column sep=small,row sep=normal,
execute at end picture={\draw (brNW.north west) rectangle (brSE.south east);
\draw (brNW.north east) -- (brSW.south west) (brNE.north east) -- (brSE.south west);}]
|[alias=brNW]| \deleted{\mathrm{Br}(\bar Y)}\arrow[r]\arrow[d,"\handwritten{\alpha}"] \&
|[alias=brNE]| \deleted{\mathrm{Br}(\bar X)}\arrow[r,"\handwritten{0}" description]\arrow[d,"\handwritten{\beta}"] \&
\deleted{H^1(\bar Y,\bar P)}\arrow[r]\arrow[d,"\handwritten{\gamma}"] \&
\deleted{H^{\handwritten3}(\bar Y,\underline G_m)}\arrow[r]\arrow[d,"\handwritten{\alpha'}"] \&
\deleted{H^{\handwritten3}(\bar X,\underline G_m)}\arrow[d,"\handwritten{\beta'}"]\\
|[alias=brSW]| \deleted{\mathrm{Br}(y)}\arrow[r] \& |[alias=brSE]| \deleted{\mathrm{Br}(X_y)}\arrow[r,"\handwritten{0}" description] \&
\deleted{H^1(y,P_{\unsure{o|0}})}\arrow[r] \& \deleted{H^{\handwritten3}(y,\underline G_m)}\arrow[r] \&
\deleted{H^{\handwritten3}(X_y,\underline G_m)}
\end{tikzcd}}
\deleted{\deleted{\handwritten{\unsure{?}}}\quad ).}

\deleted{Les flèches $\handwritten{\alpha,\beta,\alpha'}$ du diagramme sont bijectives, en vertu de (BR I\phantom{0000}), (\handwritten{3}.1), (\phantom{000}), d'où résulte que \deleted{il en est de même} $\handwritten{\gamma}:H^1(Y,P)\handwritten{\longrightarrow} H^1(\handwritten{y\deleted{\unsure{{}_y}}},P_o)$ est injective, où $P_o=\underline{\mathrm{Pic}}_{X_y/y}$. Donc (4.1\handwritten{5}) est une conséquence de $H^1(y,P_o)=0$, elle-\deleted{\unsure{m}}même conséquence d'après la deuxième suite exacte ci-dessus, des relations (4.1\handwritten{6})).}

Alors l'image de $H^1(Y,P)$ dans $H^1(Y,\underline B)$ est le sous-groupe de $H^1(Y,\underline B)$ formé des \unsure{élements|élemtns|éléments} dont l'image dans chaque $H^1(\bar Y,\underline B)$ est nulle. Ce sous-groupe s'appelle le \uline{groupe de Tate-Chafarévitch de} $Y$ \uline{à coéfficients dans} $\underline B$, et se note $\handwritten{\overline{\underline{\mathrm{III}}}}(Y,\underline B)$ (cf plus bas \inserted{pour} la relation avec la définition originelle de Tate et Chafarévitch). On trouve alors la suite exacte déduite de (4.1\deleted{\unsure{2|4}}\handwritten{4}) :
\[
\ldots\ \mathrm{Pic}(X_{\handwritten{\eta}}/\handwritten{\eta})\handwritten{\longrightarrow\coprod}_{y}H^1(y,\underline E_y)
\handwritten{\longrightarrow} H^1(Y,P)\handwritten{\longrightarrow\overline{\underline{\mathrm{III}}}}(Y,\underline B)\handwritten{\longrightarrow}0\ ,
\tag*{(4.1\deleted{\unsure{5|6}}\handwritten{7})}
\]
\margin{\handwritten{\textcircled{P}}}
explicitant $H^1(Y,\deleted{\unsure{\underline B}}\handwritten P)$ comme une extension du groupe de Tate-Chafarévitch par un quotient de la somme $\handwritten{\coprod}_yH^1(y,\underline E_y)$.

\margin{\handwritten{\unsure{PP|//}}}
Il faut enfin expliciter cette dernière somme. Notons que $\underline E_y$ \deleted{n'est non nul qu} est nul si $X_y$ est \inserted{\deleted{géométriquement}\handwritten{géométriquement}} intègre, à fortiori si $X_y$ est lisse sur $y$, de sorte que si $X_{\handwritten{\eta}}$ est séparable sur $k(\handwritten{\eta})$ (par exemple, lisse sur $k(\handwritten{\eta})$), alors le support de $E$ est un \deleted{ensemble}\inserted{contenu dans (en fait, égal à) l'ensemble} \uline{fini} des $y_i\handwritten{\in}Y$ tels que $X_{y_i}$ ne soit pas géométriquement intègre. Un calcul facile donne pour $\underline E_y$ la structure suivante (en procédant par descente à partir de $\underline E_{\bar y}$). On écrit \handwritten{l'égalité de \deleted{\unsure{cycles}}\inserted{diviseurs sur} $X$ :}
\[
X_y=\handwritten{\sum}_{i} a_y^i\,C_y^i\ ,
\tag*{(4;1\deleted{\unsure{6}}\handwritten8)}
\]
les $C_y^i$ étant les composantes irréductibles de $X_y$, et\inserted{\handwritten{ les}} $a_y^i$ leur multiplicité dans la fibre. On désigne par
% archive sheet number: 11
